\chapter{La caravane}
\label{chap:la-caravane}

Cela faisait une semaine qu'ils avaient quitté le village lorsque la forêt
commença enfin à s'éclaircir.

Depuis plusieurs heures, Khaer Ghal avait remarqué que le terrain changeait.
Les arbres s'espaçaient davantage et, sous leurs pieds, les traces d'un ancien
passage étaient devenues progressivement plus nombreuses. Ce n'était pas
encore une véritable route, mais suffisamment pour lui faire comprendre que
d'autres personnes empruntaient régulièrement cet endroit.

Kamatsu fut le premier à entendre les chevaux.

Il s'immobilisa.

— T'as entendu ?

Khaer Ghal tendit l'oreille.

Au loin, derrière les arbres, quelque chose grinçait régulièrement. Une roue,
peut-être. Puis une voix leur parvint, trop lointaine pour qu'ils puissent
comprendre les mots.

Kamatsu regarda son ami.

Depuis une semaine, ils n'avaient rencontré personne.

Ils avancèrent plus prudemment et finirent par apercevoir les premiers chariots
entre les troncs. Plusieurs véhicules étaient arrêtés le long d'un chemin plus
large, entourés de chevaux, d'hommes et de femmes occupés à préparer ce qui
semblait être une halte. Des caisses avaient été descendues, quelqu'un
s'occupait d'un attelage et une fine fumée commençait à monter au-dessus d'un
feu.

Après sept jours passés seuls dans la forêt, cette simple agitation avait
quelque chose de presque rassurant.

Kamatsu sourit.

— Des voyageurs.

Khaer Ghal ne répondit pas immédiatement.

Il observait.

Ils étaient nombreux. Beaucoup plus nombreux qu'eux, évidemment. Plusieurs
portaient des armes et au moins l'un des hommes semblait surveiller les
alentours plutôt que participer à l'installation du campement.

Pourtant, rien ne ressemblait au village qu'ils avaient quitté.

Il n'y avait ni palissade, ni guerriers orcs, ni prisonniers.

Seulement des voyageurs.

Ils avaient besoin de nourriture, de repos et surtout de savoir où ils se
trouvaient. Continuer seuls sans même connaître la direction de la prochaine
ville n'avait guère de sens.

Les deux garçons quittèrent donc les arbres.

Ils n'avaient parcouru qu'une partie de la distance lorsque l'homme qui montait
la garde les aperçut.

Son regard se posa d'abord sur Kamatsu.

Puis sur Khaer Ghal.

Il se figea.

— Un Orc !

Le cri traversa le campement.

En quelques secondes, les conversations s'interrompirent. Des hommes et des
femmes surgirent entre les chariots tandis que d'autres abandonnaient ce qu'ils
étaient en train de faire. Certains avaient saisi une épée ou un arc ; d'autres
tenaient simplement l'outil qui se trouvait à portée de leur main.

Khaer Ghal s'arrêta immédiatement.

Kamatsu aussi.

La réaction leur parut d'abord incompréhensible.

Puis Khaer Ghal vit les regards.

Ils passaient constamment de lui à Kamatsu.

Les caravaniers ne voyaient pas deux enfants qui venaient de passer une semaine
seuls dans les bois. Ils voyaient un jeune humain amaigri, sale, couvert
d'éraflures et vêtu de vêtements abîmés par plusieurs jours de marche.

À côté de lui se tenait un Orc.

Et des deux garçons, Khaer Ghal était le seul à porter une arme.

— Il a un enfant avec lui !

Plusieurs armes se tournèrent aussitôt vers lui.

Khaer Ghal recula instinctivement d'un pas.

Sa main descendit presque vers son épée avant qu'il ne s'arrête. Il comprit
immédiatement ce que ce geste provoquerait et écarta au contraire légèrement
les bras pour montrer qu'il ne comptait pas dégainer.

Kamatsu regardait autour de lui.

D'autres personnes arrivaient derrière eux, attirées par les cris. Elles ne
cherchaient pas encore à les encercler, mais retourner dans la forêt
signifierait désormais passer devant plusieurs adultes armés.

Un homme s'avança de quelques pas.

Son épée resta dirigée vers le sol, mais sa main ne quitta pas la poignée.

— Laisse partir le garçon.

Khaer Ghal fronça les sourcils.

— Quoi ?

— Le garçon. Laisse-le partir.

L'homme regarda ensuite Kamatsu et son expression s'adoucit légèrement.

— Viens par ici, petit. Ça va aller.

Kamatsu regarda le caravanier, puis Khaer Ghal, puis de nouveau le caravanier.

Il comprit.

— Mais il me retient pas !

Quelques regards s'échangèrent parmi les adultes.

Personne n'abaissa son arme.

— Viens avec nous, répéta l'homme. Il ne te fera rien.

— Mais c'est mon ami !

Cette fois, Kamatsu fit un pas.

Pas vers les caravaniers, vers Khaer Ghal. Il se plaça devant lui.

Le mouvement fut si naturel qu'il sembla n'y avoir réfléchi qu'une fois qu'il
était déjà là. Il se tenait légèrement décalé devant son ami, suffisamment pour
que son propre corps se trouve désormais entre Khaer Ghal et plusieurs des
armes qui le menaçaient.

Alors seulement lui revinrent les derniers mots que la mère de Khaer Ghal lui
avait adressés avant leur départ. Elle lui avait demandé de veiller sur son
fils, comme celui-ci avait veillé sur lui.

Kamatsu n'avait pas oublié.

— Kamatsu...

— Non.

Le jeune humain leva les deux mains devant lui, paumes ouvertes.

— Il m'a rien fait !

Khaer Ghal posa doucement une main sur son épaule.

— On peut repartir.

Kamatsu ne bougea pas.

— Ils vont te tuer si tu bouges.

\scenebreak

Khaer Ghal regarda les adultes qui leur faisaient face.

Il connaissait ce genre de peur.

Il l'avait vue chez les prisonniers du village lorsqu'un guerrier s'approchait
d'eux. Il l'avait vue dans les yeux de Kamatsu et de ses parents. Elle n'était
pas tout à fait la même ici, parce que ceux qui avaient peur étaient également
ceux qui tenaient les armes, mais il la reconnaissait malgré tout.

— Ils ont peur, dit-il.

Kamatsu essaya alors de leur expliquer.

Les mots sortirent trop vite, dans le désordre. Il parla du village, de ses
parents, de l'arène, revint à Khaer Ghal, repartit sur leur première tentative
de fuite avant de se rappeler qu'il n'avait pas expliqué pourquoi ils avaient
voulu s'enfuir.

Ce n'était pas un récit très clair. C'était celui d'un enfant fatigué qui
essayait de condenser en quelques minutes tout ce qu'ils venaient de vivre.

Lorsqu'il eut terminé, Kamatsu reprit son souffle.

— Alors arrêtez de le regarder comme ça. C'est mon ami.

Le silence qui suivit n'était plus tout à fait le même.

Les armes n'avaient pas encore disparu, mais certaines s'étaient abaissées.
Les caravaniers échangeaient désormais des regards moins assurés. Ils avaient
cru comprendre ce qu'ils avaient devant eux dès le premier instant.

Ils commençaient à envisager qu'ils s'étaient peut-être trompés.

Khaer Ghal recula lentement, en prenant soin de ne jamais approcher sa main de
son épée.

C'est à cet instant qu'un chat roux apparut entre deux chariots.

Personne ne lui prêta attention.

L'animal traversa tranquillement le groupe d'adultes, passa entre les jambes de
l'un d'eux et poursuivit son chemin comme si les armes, les cris et les deux
étrangers au milieu du campement ne constituaient qu'une perturbation
insignifiante de sa journée.

Une femme le remarqua.

— Bartiméus ?

Le chat ne lui accorda pas le moindre regard.

Il passa près de Kamatsu, le renifla brièvement, puis continua jusqu'à Khaer
Ghal.

L'Orc s'immobilisa.

Bartiméus passa entre ses jambes et se frotta contre lui.

Khaer Ghal baissa les yeux, surpris.

Le chat fit encore quelques pas avant de s'arrêter et de tourner la tête vers
lui.

Leurs regards se croisèrent.

Khaer Ghal oublia presque un instant les adultes qui les entouraient. Il
n'avait jamais vu un chat d'aussi près. Bartiméus, de son côté, semblait
l'observer avec le sérieux de quelqu'un qui venait de découvrir quelque chose
d'intéressant et comptait bien prendre le temps de se faire une opinion.

— Bartiméus...

La même femme avait parlé, mais cette fois sa voix était différente.

Un homme près d'elle baissa légèrement son arc.

— Quoi ?

— Il l'aime bien.

— C'est un chat.

— Je sais ce que c'est.

Bartiméus se frotta une nouvelle fois contre la jambe de Khaer Ghal.

La femme observa la scène quelques secondes.

— Il ne fait pas ça avec tout le monde.

Un homme eut un petit rire nerveux.

— On va vraiment faire confiance à un chat ?

— Non. Mais il se trompe rarement sur les gens.

Ce n'était certainement pas suffisant pour effacer la méfiance des
caravaniers, mais après le récit de Kamatsu, le comportement du chat semblait
avoir ébranlé un peu plus leurs certitudes.

Les armes s'abaissèrent progressivement, même si certaines restèrent encore
prudemment levées. Ce fut suffisant pour que Kamatsu baisse enfin les mains.

%\illustration
%  {0900.png}
%  {La rencontre avec la caravane}
%  {la-caravane}

La peur n'avait pas disparu et la méfiance non plus. Khaer Ghal le voyait
encore dans plusieurs regards. Mais quelque chose avait changé. Pour la
première fois depuis leur arrivée, les caravaniers semblaient réellement les
regarder.

Deux garçons, sales, épuisés, affamés, et probablement beaucoup trop jeunes
pour se trouver seuls au milieu de cette forêt.

Après quelques échanges à voix basse, les voyageurs acceptèrent finalement de
les laisser rester auprès d'eux pour la nuit. On leur donna à manger, de l'eau
et un endroit où se reposer près des chariots.

Rien de plus ne leur fut promis.

Certains continuaient à surveiller Khaer Ghal avec méfiance, tandis que
d'autres semblaient surtout s'interroger sur ce qu'ils allaient bien pouvoir
faire de ces deux enfants lorsque la caravane reprendrait la route.

Pour Khaer Ghal et Kamatsu, cela pouvait attendre.

Après une semaine passée seuls dans la forêt, ils avaient mangé à leur faim et
allaient pouvoir dormir sans avoir à chercher un abri ni à se demander ce qui
rôdait dans l'obscurité.

Pour cette nuit, c'était largement suffisant.